\documentclass[11pt,a4paper]{article}
\usepackage[utf8]{inputenc}
\usepackage[vietnamese]{babel}
\usepackage{amsmath,amssymb}
\usepackage{geometry}
\geometry{a4paper, total={180mm,260mm}, left=15mm, right=15mm, top=18mm, bottom=18mm}
\usepackage{tikz}
\usetikzlibrary{shapes.geometric, arrows.meta, positioning, calc}
\usepackage{tcolorbox}
\tcbuselibrary{skins,breakable}
\usepackage{xcolor}
\usepackage{multicol}
\usepackage{array}
\usepackage{fontawesome}
\usepackage{fancyhdr}

\pagestyle{fancy}
\fancyhf{}
\renewcommand{\headrulewidth}{0.5pt}
\renewcommand{\footrulewidth}{0.5pt}
\fancyhead[L]{\small\bfseries\color{teal}\faBook\ SỔ TAY KIẾN THỨC TOÁN 3}
\fancyhead[R]{\small\bfseries\color{teal}Lớp Toán Thầy Lĩnh -- loptoanthaylinh.vn}
\fancyfoot[C]{\small\color{gray}-- Sổ tay kiến thức: Bài 62. Luyện tập chung: Các số đến 100 000 --}

\definecolor{primary}{HTML}{005F73}
\definecolor{secondary}{HTML}{0A9396}
\definecolor{accent}{HTML}{EE9B00}
\definecolor{danger}{HTML}{AE2012}
\definecolor{boxbg}{HTML}{F4F6F8}

\begin{document}

% --- BANNER TIÊU ĐỀ ---
\begin{tcolorbox}[
  enhanced,
  colback=white,
  colframe=primary,
  arc=3mm,
  boxrule=1.2pt,
  left=4mm, right=4mm, top=2.5mm, bottom=2.5mm,
  drop shadow=black!10!white
]
  \centering
  {\Large\bfseries\color{primary}\faBookmark\ BÀI 62: LUYỆN TẬP CHUNG: CÁC SỐ ĐẾN 100 000}\\[1.5mm]
  {\small\itshape\color{black!75}Đại tổng kết cấu tạo số có 5 chữ số, so sánh dãy số, kĩ năng làm tròn số và ứng dụng đời sống trên 1 trang A4}
\end{tcolorbox}

\vspace{-2mm}


\noindent
\begin{minipage}[t]{0.49\textwidth}\vspace{0pt}
\begin{tcolorbox}[
  enhanced,
  colback=blue!4!white,
  colframe=blue!75!black,
  arc=2mm, boxrule=0.8pt,
  left=2mm, right=2mm, top=1.8mm, bottom=1.8mm,
  title={\footnotesize\bfseries\color{white}\faKey\ 1. Ma trận kiến thức số đến 100 000}
]
{\footnotesize
\begin{itemize}\setlength{\itemsep}{1.5pt}\setlength{\leftskip}{-2.5mm}
  \item \textbf{Cấu tạo 5 hàng:} Chục nghìn, nghìn, trăm, chục, đơn vị.
  \item \textbf{So sánh:} Số nào nhiều chữ số hơn thì lớn hơn; cùng số chữ số thì so sánh từ hàng cao nhất (chục nghìn).
  \item \textbf{Làm tròn 4 cấp độ:}
  \begin{itemize}\setlength{\itemsep}{0.5pt}\setlength{\leftskip}{-2.5mm}
    \item Tròn chục: xét Đơn vị.
    \item Tròn trăm: xét Chục.
    \item Tròn nghìn: xét Trăm.
    \item Tròn chục nghìn: xét Nghìn.
  \end{itemize}
\end{itemize}
}
\end{tcolorbox}
\end{minipage}
\hfill
\begin{minipage}[t]{0.49\textwidth}\vspace{0pt}
\begin{tcolorbox}[
  enhanced,
  colback=teal!4!white,
  colframe=teal!75!black,
  arc=2mm, boxrule=0.8pt,
  left=2mm, right=2mm, top=1.8mm, bottom=1.8mm,
  title={\footnotesize\bfseries\color{white}\faSliders\ 2. Kĩ năng phân tích và lập số}
]
{\footnotesize
\begin{itemize}\setlength{\itemsep}{1.5pt}\setlength{\leftskip}{-2.5mm}
  \item \textbf{Viết số thành tổng các hàng:}
  $$47\,256 = 40\,000 + 7000 + 200 + 50 + 6$$
  \item \textbf{Số tròn chục nghìn:} Có tận cùng là 4 chữ số $0$ (ví dụ $10\,000, 20\,000, \dots, 90\,000$).
  \item \textbf{Dãy số cách đều:} Tìm khoảng cách giữa 2 số liền kề rồi cộng tiếp để điền số còn thiếu.
\end{itemize}
}
\end{tcolorbox}
\end{minipage}

\vspace{-2.5mm}
% --- KHỐI 3: SƠ ĐỒ TƯ DUY TIKZ ---
\noindent
\begin{tcolorbox}[
  enhanced,
  colback=orange!4!white,
  colframe=orange!80!black,
  arc=2mm, boxrule=0.8pt,
  left=2mm, right=2mm, top=1.5mm, bottom=1.5mm,
  title={\footnotesize\bfseries\color{white}\faSitemap\ SƠ ĐỒ TƯ DUY: HỆ SINH THÁI CÁC SỐ ĐẾN 100 000}
]
\centering
\begin{tikzpicture}[
  box/.style={draw, rounded corners=2pt, align=center, font=\scriptsize\bfseries, inner sep=2.5pt, minimum height=0.75cm},
  arr/.style={-{Stealth[length=2mm]}, thick}
]
  \node[box, fill=purple!15, draw=purple!80!black, text width=2.6cm] (m1) at (0, 0) {ĐỌC VIẾT SỐ\\5 chữ số thập phân\\$100\,000$ có 6 chữ số};
  \node[box, fill=teal!15, draw=teal!80!black, text width=2.6cm] (m2) at (3.5, 0) {SO SÁNH THỨ TỰ\\Xét từ trái sang phải\\Tăng / Giảm dần};
  \node[box, fill=blue!12, draw=blue!80!black, text width=2.6cm] (m3) at (7.0, 0) {LÀM TRÒN SỐ\\Quy tắc mốc 5\\Tròn nghìn, chục nghìn};
  \node[box, fill=orange!15, draw=orange!80!black, text width=2.6cm] (m4) at (10.5, 0) {ỨNG DỤNG\\Dân số, độ dài\\Giá tiền sản phẩm};

  \draw[arr, purple!80!black] (m1) -- (m2);
  \draw[arr, teal!80!black] (m2) -- (m3);
  \draw[arr, blue!80!black] (m3) -- (m4);
\end{tikzpicture}
\end{tcolorbox}

\vspace{-2.5mm}
% --- KHỐI 4: VÍ DỤ & MẸO TRÁNH BẪY ---
\noindent
\begin{minipage}[t]{0.49\textwidth}\vspace{0pt}
\begin{tcolorbox}[
  enhanced,
  colback=green!4!white,
  colframe=green!60!black,
  arc=2mm, boxrule=0.8pt,
  left=2mm, right=2mm, top=1.8mm, bottom=1.8mm,
  title={\footnotesize\bfseries\color{white}\faCheckCircle\ Ví dụ thực hành chuẩn mực}
]
{\footnotesize
\textbf{1. Lập số từ 5 chữ số $\{2; 0; 7; 4; 9\}$:}\\
-- Số lớn nhất: $\mathbf{97\,420}$.\\
-- Số bé nhất: $\mathbf{20\,479}$ (hàng chục nghìn phải khác $0$).\\
\textbf{2. Ứng dụng thực tế:}\\
Đoạn đường dài $48\,620\text{ m}$.\\
-- Làm tròn đến hàng nghìn: $\mathbf{49\,000\text{ m}}$.\\
-- Làm tròn đến hàng chục nghìn: $\mathbf{50\,000\text{ m}}$.
}
\end{tcolorbox}
\end{minipage}
\hfill
\begin{minipage}[t]{0.49\textwidth}\vspace{0pt}
\begin{tcolorbox}[
  enhanced,
  colback=red!4!white,
  colframe=red!70!black,
  arc=2mm, boxrule=0.8pt,
  left=2mm, right=2mm, top=1.8mm, bottom=1.8mm,
  title={\footnotesize\bfseries\color{white}\faExclamationTriangle\ Mẹo tránh bẫy trắc nghiệm}
]
{\footnotesize
\begin{itemize}\setlength{\itemsep}{1.5pt}\setlength{\leftskip}{-2.5mm}
  \item \textbf{Bẫy tách số có chữ số 0:}\\
  $50\,402 = 50\,000 + 400 + 2$. Nhiều bạn viết thừa $+ 0000 + 00$ gây rối mắt, hoặc viết thiếu số $0$ thành $5000 + 400 + 2$ là sai!
  \item Khi đọc số tiền thực tế, luôn phân biệt giữa "chục nghìn" ($10\,000$) và "trăm nghìn" ($100\,000$).
\end{itemize}
}
\end{tcolorbox}
\end{minipage}


\end{document}
